% !TeX root = ../main.tex

\hfutsetup{
  keywords  = {遥感图像, 微小船舶检测, LEVIR-Ship, 目标检测, 检测系统},
  keywords* = {remote sensing images, tiny ship detection, LEVIR-Ship, object detection, detection system},
}

\begin{abstract}
  遥感场景下的微小船舶检测在海上监管、港口监测、海事执法和应急救援等应用中具有重要价值。然而，中分辨率遥感图像中的船舶目标往往尺度小、纹理弱，并且易受到海浪、云层、岸线和港区设施等复杂背景干扰，导致检测任务普遍存在漏检率高、误检率高以及模型跨场景稳定性不足等问题。围绕上述问题，本文以公开数据集 LEVIR-Ship 为统一实验基准，开展了多模型对比研究与检测系统实现工作。

  在算法研究部分，本文首先梳理了通用目标检测、小目标检测和遥感船舶检测相关研究，在统一数据划分、评价指标和结果记录方式的基础上，选取 DRENet、FCOS 和 YOLO26 三类代表模型进行对比实验。实验采用 AP50 作为主指标，并结合 AP50-95、Precision、Recall 和 F1 对模型性能进行综合评价。当前实验结果表明，DRENet 在 LEVIR-Ship 数据集上的 AP50 为 0.7949、Recall 为 0.8511，表现出较强的目标检出能力；YOLO26 的 AP50 为 0.7950、AP50-95 为 0.3170、Precision 为 0.8430，在精度、定位质量和工程可部署性之间表现较为均衡；FCOS 的 AP50 为 0.7700，能够作为 anchor-free 路线的参考基线纳入比较。综合来看，DRENet 更适合强调漏检控制的应用场景，YOLO26 更适合作为轻量部署候选，FCOS 则为不同检测范式的比较提供了补充视角。

  在系统实现部分，本文将现有推理能力封装为一套面向本科毕业设计场景的轻量级检测系统。系统采用 React 构建前端页面，采用 FastAPI 构建后端服务，并通过 SQLite 完成任务、结果和产物路径的持久化存储。系统支持单图与批量任务提交、模型启停、异步任务调度、检测框可视化、结构化结果回看和历史任务查询，形成了从模型推理到结果展示的完整链路。相关实现表明，该系统能够满足毕业设计阶段对“可运行、可演示、可追溯”的基本要求。

  本文工作表明，在统一评测协议下开展多模型对比，并进一步完成阈值消融、输入尺寸敏感性分析以及算法服务化和可视化系统实现，能够更全面地评估不同检测范式在遥感微小船舶场景中的适用性。本文形成的实验结果、分析框架和系统实现，可为后续模型优化、融合评测和工程部署研究提供基础。
\end{abstract}

\begin{abstract*}
  \noindent
  Tiny ship detection in remote sensing images is important for maritime surveillance, port monitoring, law enforcement and emergency response. However, ship targets in mid-resolution remote sensing images are usually very small, weak in texture, and easily affected by complex backgrounds such as waves, clouds, coastlines and harbor facilities. These characteristics make the task prone to missed detections, false alarms and unstable cross-scene performance. To address this problem, this thesis conducts comparative experiments and system implementation based on the public LEVIR-Ship dataset.

  In the algorithm study, related work on general object detection, small object detection and remote sensing ship detection is first reviewed. Then, under a unified data split, evaluation protocol and result recording scheme, three representative models, namely DRENet, FCOS and YOLO26, are selected for comparative experiments. AP50 is used as the primary metric, while AP50-95, Precision, Recall and F1 are used as complementary metrics. The current results show that DRENet achieves an AP50 of 0.7949 and a Recall of 0.8511 on LEVIR-Ship, indicating strong target discovery capability. YOLO26 achieves an AP50 of 0.7950, an AP50-95 of 0.3170 and a Precision of 0.8430, showing a balanced performance in terms of accuracy, localization quality and deployability. FCOS achieves an AP50 of 0.7700 and is retained as an anchor-free reference baseline. Overall, DRENet is more suitable for scenarios that prioritize recall, while YOLO26 is more suitable for lightweight deployment.

  In the system implementation part, the existing inference capability is encapsulated into a lightweight detection system tailored to the undergraduate thesis scenario. The frontend is built with React, the backend is built with FastAPI, and SQLite is used to persist tasks, results and generated artifacts. The system supports single-image and batch task submission, model management, asynchronous execution, visualization of detection boxes, structured result review and history query, thereby forming a complete workflow from inference to presentation. The implementation demonstrates that the system can satisfy the basic requirements of being executable, demonstrable and traceable.

  The study shows that combining unified multi-model comparison with threshold ablation, input-size sensitivity analysis, and service-oriented system implementation can provide a more comprehensive understanding of the applicability of different detection paradigms to tiny ship detection in remote sensing scenes. The experimental results, analysis framework and system implementation can provide a foundation for subsequent model optimization, fusion evaluation and engineering deployment.
\end{abstract*}
